\scanpage{111}
\mspage{345}

\deleted{et on a en fait alors
\[
g\rho(g)\rho^2(g)=g^3=\rho^{-3}=-1
\]}
\inserted{$=\begin{pmatrix}0&\unsure{-1|+1}\\1&-1\end{pmatrix}$}

\inserted{Remarques} \deleted{\unsure{NB}} 1°) Evidemment $-\rho$ est aussi solution de (\unsure{87})
(tout comme $\rho^{-1}$) \deleted{\unsure{et $\rho$ donc} $(-\rho)\rho(g)\rho^2(g)=1$}
(tout comme $-\rho^{-1}$), mais il est justiciable
respectivement des cas \deleted{\illegible{1}} 1°) \deleted{\unsure{et du cas}}
\deleted{\unsure{A, resp.}}

2°) Pour \deleted{\illegible{1}}\inserted{déterminer} les solutions dans $\mathrm{Gl}(2,\widehat{\mathbb Z})$
de $g\rho(g)\rho^2(g)=1$, il suffit de choisir pour
\unsure{$\forall\,\ell$} une telle solution dans $\mathrm{Gl}(2,\widehat{\mathbb Z})$, en
utilisant les \unsure{résultat|résultats} (86) ; on notera que
les ``cas'' obtenus 1°) à 3°) peuvent varier d'un
$\ell$ à un autre. \deleted{Même remarque pour
les solutions de l'équation $g\rho(g)\rho^2(g)=\unsure{-1|1}$.}
On fera attention que les cas 1°), 2°), 3°)
correspondent \unsure{resp.\textsuperscript{t}} à des familles de
solutions dépendant de 2, 1 et 0 paramètres.

\unsure{3°)} J'ignore si une $\xi\in\mathrm{Sl}(2,\widehat{\mathbb Z})^{\unsure{\prime|*}}$, \inserted{$\xi=\begin{pmatrix}A&B\\C&D\end{pmatrix}$}, \inserted{qui satisfait}
$g\rho(g)\rho^2(g)=1$ admet une représentation
paramétrique de la forme (58)
par $\lambda_0\rho(\lambda_0)^{-1}$ ou $\lambda'_0(-l_1^{-1})\rho(\lambda'_0)^{-1}$
\inserted{$=\lambda_0(-\unsure{l_1^{-1}|l_0^{-1}})\rho(\lambda_0)^{-1}$ avec $\lambda_0\unsure{=\lambda'_0\tau_\infty}$}
(il suffit
de regarder séparément dans les $\mathrm{Sl}(2,\mathbb Z_\ell)'$).
Je l'ignore $\overline{\mathrm m}$ si on suppose $g\equiv\mathrm{id}\ (2)$,
(ce qui est une condition sur la composante
dans $\mathrm{Sl}(2,\mathbb Z_2)'$ seulement). On voit en

\margin{\uline{NB} $g=l_0^{-1}$ n'est \uline{pas}
solution de cette \unsure{l'équation|équation}, par contre
$g=-l_0^{-1}$ l'est\ldots}

\scanpage{112}

\unsure{tous|tout} cas sur les formules (78) que \unsure{\uline{si}} $\xi$
\deleted{condition}\inserted{est de} la forme $\lambda_0\rho(\lambda_0)^{-1}$
\inserted{(si $\xi\rho(\xi)\rho^2(\xi)=1$, bien sûr)},
\deleted{\unsure{a bien alors le cas A)}} \deleted{\illegible{1}}
\inserted{on aura} $C+D-B=1$ --- i.e. \deleted{\illegible{1}}\inserted{on est}
\inserted{dans le} cas ``général'' à deux paramètres.
\deleted{Je présume
que si on prend
\[
g=\lambda_0(-l_0^{-1})\rho(\lambda_0)^{-1},
\]
on sera dans le cas B), avec $c+d-b=-1$
i.e. le cas 2°) encore à deux paramètres. Je
\unsure{me dispense} de faire le calcul ici.}

\margin{OK, cf. (\unsure{99}) ci-dessous.}

(Sans pouvoir néanmoins exclure le cas
\inserted{où on devrait en $\overline{\mathrm m}$ temps}
$A+B=0$, \deleted{ce qui signifie en fait que
\[
\xi=\begin{pmatrix}A&B\\C&D\end{pmatrix}
=-\rho^{-1}=\begin{pmatrix}-1&1\\-1&0\end{pmatrix}
\]}
i.e. le cas à \unsure{un} paramètre
\[
\xi=\begin{pmatrix}-u&u\\C&-u'-C\end{pmatrix}
\]
avec $u$ racine primitive \unsure{3ième} de 1,
$u'=u^{-1}=u^2$ l'autre racine primitive.

\scanpage{113}
\mspage{346}

(VIII) --- \uline{Calculs dans} $\mathrm{Sl}(2,\widehat{\mathbb Z})$ : \uline{cocycles sous} $\sigma_\infty$.

Reprenons un $g\in\begin{pmatrix}a&b\\c&d\end{pmatrix}\in\mathrm{Sl}(2,\widehat{\mathbb Z})$,
et calculons $\sigma_\infty g\sigma_\infty^{-1}$, on trouve, grâce à
\[
\sigma_\infty=\begin{pmatrix}0&-1\\1&0\end{pmatrix}\qquad (48)
\]
\deleted{\illegible{1}}
les formules
\[
(90)\quad\left\{
\begin{aligned}
\sigma_\infty g\sigma_\infty^{-1}
&=\begin{pmatrix}d&-c\\-b&a\end{pmatrix}={}^t(g^{-1})
&&\text{la contragrédiente de $g$ !}\\
&\text{i.e.}\\
\sigma_\infty g^{-1}\sigma_\infty^{\unsure{-1}}
&={}^t g&&\text{donc}\\
[\sigma_\infty,g^{-1}]&=({}^t g)(g)
=\begin{pmatrix}a^2+b^2&ac+bd\\ac+bd&\unsure{c^2+d^2|c^2-d^2}\end{pmatrix}
\end{aligned}\right.
\]
\inserted{$=\sigma_\infty g^{-1}\unsure{\sigma_\infty|\sigma_\infty^{-1}}g$}

\margin{\uline{NB} on a
$\sigma_\infty^{-1}=\begin{pmatrix}0&1\\-1&0\end{pmatrix}=-\sigma_\infty$
et \deleted{\illegible{1}} $\sigma_\infty^2\ne1$ dans $\mathrm{Sl}(2,\widehat{\mathbb Z})$ !}

\margin{\uline{NB} Si $g\in\mathrm{Gl}(2,A)$,
$A$ anneau quelconque
(p. ex. $A=\mathbb Z_\ell$), on a
$\sigma_\infty g\sigma_\infty^{-1}=\begin{pmatrix}d&-c\\-b&a\end{pmatrix}
=(\det g){}^t(g^{-1})$}

donc, en prenant l'inverse des deux mbs dans
la dernière formule
\[
(91)\quad \underbrace{g^{-1}\sigma_\infty(g)}_{\inserted{({}^t g\cdot g)^{-1}}}
=\begin{pmatrix}c^2+d^2&-(ac+bd)\\-(ac+bd)&a^2+b^2\end{pmatrix}.
\]

D'autre part, soit
\deleted{$h_0=\begin{pmatrix}\unsure{r}&s\\t&u\end{pmatrix}$}
$h_0\in\mathrm{Sl}(2,\widehat{\mathbb Z})$
on a
\[
h_0\sigma_\infty(h_0)=h_0({}^t h_0^{-1})
\]
donc
\[
(*)\qquad h_0\sigma_\infty(h_0)=1\quad\Longleftrightarrow\quad h_0={}^t h_0.
\]

Soit enfin
\[
(91)\qquad \xi=\begin{pmatrix}A&B\\C&D\end{pmatrix}
\]
et considérons l'équation (68)
\[
(**)\qquad \sigma_\infty(\xi)=l_1^\gamma\xi^{-1}l_0^\gamma
\]
qui équivaut aussi à (*), avec $h_0$ donné par (\unsure{65|66})
\deleted{$h_0=l_0^{-\unsure{\gamma|\delta}}\xi$ \illegible{4}}, $l_0,l_1$ par (44).

\scanpage{114}

\[
(92)\quad\begin{aligned}
h_0=l_0^{-\gamma}\xi
&=\begin{pmatrix}1&+2\gamma\\0&1\end{pmatrix}
\begin{pmatrix}A&B\\C&D\end{pmatrix}\\
&=\begin{pmatrix}A+2\gamma C&B+2\gamma D\\C&D\end{pmatrix}
\end{aligned}
\]
donc ``l'équation en $\xi$'' (**) (écrite dans le
quotient $\mathrm{Sl}(2,\widehat{\mathbb Z})'$ de $\widehat{\mathfrak S}_{0,3}^{+}$, s'écrit $B+2\gamma D=C$
i.e.
\[
(93)\qquad \underbrace{C-B}_{\inserted{1-D}}=2\gamma D
\qquad\text{i.e.}\quad\boxed{pD=1}\ \inserted{\text{\unsure{en cas|un cas}}}
\]
\deleted{\illegible{2}}

\margin{en se \unsure{rappelant} que $2\gamma+1=p$}

(On fera attention que l'équation (**) \uline{dans}
$\mathrm{Sl}(2,\widehat{\mathbb Z})'$ signifie $h_0\sigma_\infty(h_0)=\pm1$ i.e. $h_0=\pm{}^t h_0$,
\deleted{donc ici} mais le signe $-$ impliquerait qu'on a
$D=-D$ i.e. $2D=0$ donc $D=0$, ce qui est impossible
si on suppose que $\xi=\begin{pmatrix}A&B\\C&D\end{pmatrix}$ provient
de $\widehat\Pi_{0,3}$ donc est dans $\Pi'(2)$ i.e.
$\equiv\deleted{\unsure{\pm}}1\ (2)$ (de sorte que $D\equiv1\ (2)$).
Donc on devra bien avoir (93).

Prenons maintenant la représentation
paramétrique \inserted{(70)} de $\xi$ (tenant compte de
l'équation $\xi\rho(\xi)\rho^2(\xi)=1$ \deleted{\unsure{$\xi$}} (64))
\[
(94)\qquad\begin{aligned}
\xi&=\lambda\rho(\lambda)^{-1}
&&\inserted{(\lambda\in\mathrm{Sl}(2,\widehat{\mathbb Z}))}\\
&\text{ou}\quad
\xi=\lambda'(\unsure{-l_1^{-1}})\rho(\lambda')^{-1}
=\lambda\rho(\lambda)^{-1}
&&\inserted{\left(\begin{gathered}\lambda=\lambda'\tau_\infty\\\lambda'\in\mathrm{Sl}(2,\widehat{\mathbb Z})\end{gathered}\right)}
\end{aligned}
\]
\deleted{\unsure{$\lambda,\lambda'\in\mathrm{Sl}$}}
avec $\lambda\in\mathrm{Gl}(2,\widehat{\mathbb Z})^{\unsure{*|\circ}}$ \inserted{\deleted{\unsure{cas I}} $\det\lambda=\pm1$},
(en fait $\lambda\equiv\deleted{\unsure{\pm}}1\ (2)$ dans le
cas qui nous intéresse, mais peu importe\ldots)

\deleted{A) Cas $\xi=\lambda\rho(\lambda)^{-1}=[\lambda,\rho]$ avec $\lambda=\begin{pmatrix}a&b\\c&d\end{pmatrix}$,}

\scanpage{115}
\mspage{347}

alors la formule (78) donne $A,B,C,D$ comme
fonctions quadratiques en $a,b,c,d$, \deleted{on aura}
\inserted{(en faisant Attention \unsure{au} signe $\delta=-1$ dans le cas 2°)}
\deleted{
\[
(95)\quad\left\{
\begin{aligned}
\unsure{B-C|B+C}&=\unsure{\delta(c^2+d^2+cd-\delta)}=D-1\\
D&=\delta(c^2+d^2+cd)
\end{aligned}\right.
\]
\unsure{$D=$}}
donc (93) s'écrit
\deleted{$1-D=2\gamma D$ \unsure{ou} i.e.}
\[
(96)\qquad\boxed{(2\gamma+1)D=1\qquad\text{i.e.}\quad pD=1}
\qquad D=\delta(c^2+d^2+cd),
\]
i.e. le multiplicateur $p$ s'obtient en fonction
de $\xi\in\widehat\Pi_{0,3}$ en prenant l'image
de $\xi$ dans $\mathrm{Sl}(2,\widehat{\mathbb Z})'$, sous forme $\begin{pmatrix}A&B\\C&D\end{pmatrix}$
(il y a une ambiguïté de signe, levée par le
choix de $\lambda$ tel que $\xi=\lambda\rho(\lambda)^{-1}$, l'image
de $\lambda$ dans $\mathrm{Gl}(2,\widehat{\mathbb Z})'$ \unsure{ne|en} déterminant une
matrice (notée encore $\lambda$) de $\mathrm{Gl}(2,\widehat{\mathbb Z})$ qu'à un
\uline{signe} près, mais $\lambda\rho(\lambda)^{-1}$ ne dépend plus du
choix de ce \uline{signe}\ldots), par la formule
\[
(97)\qquad p=\frac1D\qquad (D=\unsure{\delta}(c^2+d^2+cd))
\]

Notons maintenant qu'en réduction mod 3,
la forme quadratique $c^2+d^2+cd=(c+d)^2-cd$
ne prend que les valeurs 0 et 1 (clair si $c=0$,
si $d=0$, \deleted{ou $c+d=0$ alors $f(c,d)=-cd=c^2$}
--- donc on
trouve un carré parfait dans ces \deleted{trois} cas, sinon
$c,d=\pm1$, et $f(c,d)=1+1\pm1\equiv1$ ou $0$ \deleted{\illegible{1}} (3). Or ici
on doit trouver une unité donc la valeur 0 exclue,
donc

\scanpage{116}

\margin{\uline{NB} On a $D\equiv p\equiv\delta\ (3)$ (98) la connaissance de $D$ (ou $p$, ou $\delta$) mod $3$ nous \deleted{dit \unsure{déjà|donc}} \unsure{dit} --- quel cas on a dans la \unsure{repr.} param. de $\xi$. Inutile maintenant de traiter \unsure{ces} cas séparément --- puisqu'on a $\xi=\lambda\rho(\lambda)^{-1}$, $\lambda=\lambda'\tau_\infty$, $\det\lambda=-1$.}

\[
p\equiv\delta\quad(3)
\]
\inserted{(où $\delta=\det\lambda\in\{\pm1\}$, $\xi=\lambda\rho(\lambda)^{-1}$)}
\deleted{\illegible{4} $\xi=\lambda\rho(\lambda)^{-1}$, $\det\lambda=\delta\in\{\pm1\}$, \illegible{3}}
i.e.
\[
\left\{\begin{aligned}
\text{Cas }\xi&=\lambda\rho(\lambda)^{-1}\quad(\lambda\in\mathrm{Sl}(2,\widehat{\mathbb Z}))
\Longleftrightarrow p\equiv1\ (3)\quad\deleted{\unsure{g=\lambda_0\rho(\lambda_0)^{-1}}}\Longleftrightarrow\gamma\equiv0\ (3)\\
\text{Cas }\xi&=\lambda'l_1^{-1}\unsure{\rho(\lambda)^{-1}|\rho(\lambda')^{-1}}\quad(\lambda'\in\mathrm{Sl}(2,\widehat{\mathbb Z}))
\Longleftrightarrow p\equiv-1\ (3)\Longleftrightarrow\gamma\equiv-1\ (3).
\end{aligned}\right.
\]

B) Cas $\xi=\lambda'l_0^{-1}\rho(\lambda')^{-1}=(\lambda'l_0^{-1})(\rho\lambda'^{-1}\rho^{-1})$.

Utilisant la deuxième formule (77) on trouve
\[
\xi=\begin{pmatrix}a&b\\c&d\end{pmatrix}
\begin{pmatrix}1&2\\0&1\end{pmatrix}
\begin{pmatrix}a-c&c\\a+b-c-d&c+d\end{pmatrix}
=\begin{pmatrix}A&B\\C&D\end{pmatrix}
\]
avec
\[
(99)\quad
\left\{\begin{aligned}
A&=3a^2+b^2+3ab-3ac-2ad-bc-bd
& B&=3ac+2ad+bc+bd\\
C&=-(3c^2+3cd+\deleted{\unsure{2}}d^2)+3ac+ad+2bc+\deleted{\unsure{2}}bd
& D&=3c^2+3cd+d^2
\end{aligned}\right.,
\]
\margin{\uline{NB} $C+D-B=\unsure{-1|-\delta}$.}

et la formule (93) (qui exprime ``l'équation en $\xi$'' (**), ou celle en $h_0$ (*)) s'écrit, \deleted{\unsure{donc encore sous la forme}} compte tenu de
\[
B-C=-1+D,
\]
sous la forme \deleted{\unsure{$2\gamma$}}
\[
2\gamma D=-1-D\qquad\text{i.e.}
\]
\[
(100)\quad \boxed{(2\gamma+1)D=-1\qquad\text{i.e.}\quad pD=-1},
\]
\margin{comme tantôt on ne peut avoir $h_0=-{}^t h_0$, car on ne peut avoir $D=0$ si on \unsure{suppose} \unsure{que} $\xi\equiv1\ (2)$, \unsure{donc} $D\equiv1\ (2)$.}

donc ici le multiplicateur $p$ s'exprime en fonction de l'image de $\xi$ dans $\mathrm{Sl}(2,\widehat{\mathbb Z})'$, mise sous forme $\begin{pmatrix}A&B\\C&D\end{pmatrix}$ où $A,B,C,D\in\widehat{\mathbb Z}$ --- et en levant l'ambiguïté de signe \unsure{comme tantôt} --- par la formule
\[
(101)\quad p=-\frac1D\qquad(D=\unsure{3c^2+3cd+d^2|3(c^2+3cd+d^2)}).
\]

\scanpage{117}
\mspage{348}

Ici il est clair que les seules valeurs mod $3$ \unsure{représentées} par la forme quadratique $3c^2+3cd+d^2\equiv d^2\ (3)$ sont $0$ et $1$, d'où
\[
(102)\quad p\equiv-1\ (3)\quad(\text{i.e. }\gamma\equiv-1\ (3))\quad\text{si }\xi=\lambda l_0^{-1}\rho(\lambda)^{-1}
\]
\[
(\Longleftrightarrow g=\lambda_0l_0^{-1}\rho(\lambda_0)^{-1})
\]

Cette étude montre donc bien que l'on est dans le cas $\xi=\lambda\rho(\lambda)^{-1}$, resp. $\xi=\lambda l_0^{-1}\rho(\lambda)^{-1}$, suivant que $p=\chi(u)\equiv1$ ou $\equiv-1\ (3)$, i.e. suivant que $\gamma\equiv0$ ou $\gamma\equiv-1\ (3)$.

\uline{Remarques.} 1°) Se \unsure{plaçant} dans $\mathrm{Sl}(2,\widehat{\mathbb Z})'$ \deleted{ou bien} \inserted{(ou plutôt, dans son} \deleted{du $\unsure{\Gamma|F}$} sous-groupe $\Pi(2)'=\operatorname{Ker}(\mathrm{Sl}(2,\widehat{\mathbb Z})'\unsure{-|\longrightarrow}\mathrm{Sl}(2,\mathbb Z/2\mathbb Z))$ en lieu de $\widehat{\Pi}_{0,3}$, pour \deleted{\unsure{résoudre l'équation}} \inserted{déterminer les} solutions de ``l'équation en $\xi$'' (**) (cf. (68)), quand $\xi$ est mis sous forme \unsure{paramétrique} (94), on voit qu'en tant que condition sur $\lambda=\begin{pmatrix}a&b\\c&d\end{pmatrix}$, cette équation est relativement anodine et très simple --- elle dit simplement que $pD=1$, où $D=c^2+d^2+cd$ dans le premier cas, $D=\deleted{-(3c^2+3cd+d^2)}\inserted{-(c^2+d^2+cd)}$ dans le deuxième. Si on ne fixe pas $p\in\widehat{\mathbb Z}^{*}$ d'avance, cette condition est simplement que l'élément

\scanpage{118}

$D$ de $\widehat{\mathbb Z}$ soit une unité, i.e.
\[
(103)\quad D\in\widehat{\mathbb Z}^{*}
\]
--- et on détermine alors $p$ de façon unique en termes de $D$ \inserted{, donc de $c$ et de $d$ et du signe $\delta$ par $D=\delta(c^2+d^2+cd)$, \deleted{et de $\delta\in\{\pm1\}$}} ($a$ et $b$ sont inutiles) --- donc aussi en termes de $\lambda$\inserted{$\in\widehat{\Pi}_{0,3}\cdot\unsure{(1,\tau)|\{1,\tau\}}$}, lui-$\overline{\mathrm m}$ uniquement déterminé par $\xi$.
\margin{\uline{NB} On \unsure{aura} $p$}

\uline{Ainsi la connaissance de} $\xi\in\widehat{\Pi}_{0,3}$ \uline{suffit} \uline{à déterminer} \inserted{\uline{le multiplicateur}} $p=2\gamma+1$\inserted{$\in\widehat{\mathbb Z}^{*}$}, donc aussi $\gamma\in\widehat{\mathbb Z}$, et $h_0=l_0^{-\gamma}\xi$, donc aussi $g_0$ (par la formule (66), $g_0^{-1}\sigma_\infty(g_0)=h_0$), et par suite $g$ par la formule (62), $g=\sigma_\infty(g_0)l_1^\gamma\rho(g_0)^{-1}$, donc $u$ par (50). Il est \unsure{intéressant} de s'assurer que \unsure{$D$|ça} \deleted{\illegible{1}} ne change pas quand on fait le changement (69)
\[
\xi\longmapsto\xi'=l_0^{-\alpha}\xi l_1^\alpha,
\]
en termes matriciels, on a en effet
\[
(104)\quad l_0^{-\alpha}\xi l_1^\alpha=
\begin{pmatrix}1&2\alpha\\0&1\end{pmatrix}
\begin{pmatrix}A&B\\C&D\end{pmatrix}
\begin{pmatrix}1&0\\2\alpha&1\end{pmatrix}
=\begin{pmatrix}A+2\alpha(B+C)+4\alpha^2D&B+2\alpha D\\C+2\alpha D&D\end{pmatrix}.
\]

2/ Mais on fera attention qu'en \unsure{toute} \uline{rigueur} la seule connaissance de \uline{l'image de $\xi$ dans} $\mathrm{Sl}(2,\widehat{\mathbb Z})'$ ne permet de déterminer le multiplicateur $p$ \uline{qu'au signe près} --- par contre la connaissance de l'image de $\lambda$ dans $\mathrm{Gl}(2,\widehat{\mathbb Z})'$ --- donc d'une matrice $\begin{pmatrix}a&b\\c&d\end{pmatrix}\in\unsure{\mathrm{Sl}|\mathrm{Gl}}(2,\widehat{\mathbb Z})$ définie au \uline{signe} près --- permet de déterminer

\scanpage{119}
\mspage{349}

$p$ comme l'inverse de $D$, donné \deleted{suivant les cas} par $D=\delta(c^2+cd+d^2)$ \deleted{\unsure{ou} $D=3(c^2+cd)+d^2$}. \inserted{où $\delta\in\{\pm1\}$ est déterminé par la condition $D\equiv\delta\ (3)$.}

3/ On sait par voie arithmétique (comme $\Pi_{\mathbb Q}\xrightarrow{\unsure{\chi}}\widehat{\mathbb Z}^{*}$ est surjectif) qu'on doit trouver n'importe quel $p\in\widehat{\mathbb Z}^{*}$ sous la forme $1/D$ (ou encore, sous la forme $D$), où $D=c^2+cd+d^2$ si $p\equiv1\ (3)$, et $D=-3(c^2+cd)\unsure{-|+}d^2$ si $p\equiv-1\ (3)$, avec $c,d\in\widehat{\mathbb Z}$, et $\overline{\mathrm m}$ $c\equiv0\ (2)$, $d\equiv1\ (2)$. \inserted{Si on veut la vérification directe,} (On est ramené à le vérifier séparément pour les différents $\mathbb Z_\ell$, i.e. à prouver le

\uline{Corollaire.} \unsure{1°)|°)} Si $\ell$ \deleted{\illegible{1}} premier $\ne3$, alors $\forall\unsure{\beta|p}\in\mathbb Z_\ell^{*}$ se met indifféremment sous les formes $c^2+cd+d^2$, ou $-(3(c^2+cd)+d^2)$ (avec $c,d\in\unsure{\widehat{\mathbb Z}_\ell|\mathbb Z_\ell}$). Si $\ell=2$, on peut $\overline{\mathrm m}$ supposer $c\equiv0$, $d\equiv1\ (2)$. Pour $p\in\mathbb Z_\ell^{*}$, $\ell=3$, il peut se mettre sous la forme $c^2+cd+d^2$ (resp. $-3(c^2+cd)\unsure{-|+}d^2$) ssi $p\equiv1\ (3)$ (resp. $p\equiv-1\ (3)$), cas qui sont mutuellement exclusifs.
\margin{on s'en fiche \unsure{du|de} $-(3(c^2+cd)+d^2)$ !}

4/ \unsure{Revenant} à la réduction ``anodine'' dans $\mathrm{Sl}(2,\widehat{\mathbb Z})'$ de ``l'équation en $\lambda$'' comme signifiant simplement que $c^2+d^2+cd\in\widehat{\mathbb Z}^{*}$ \deleted{resp. $3(c^2+cd)+d^2\in\widehat{\mathbb Z}^{*}$}, on pourrait espérer

\scanpage{120}

tenir une méthode pour construire $\overline{\mathrm m}$ dans $\widehat{\Pi}_{0,3}$ des gros paquets de solutions de l'équation en question. Or je n'y suis absolument pas arrivé --- pas $\overline{\mathrm m}$ d'en construire explicitement une seule, en dehors de \unsure{celles} qui correspondent au cas $u=1$, et $u=\tau_\infty$.

Une question : Soit $\lambda_0\in\mathrm{Sl}(2,\widehat{\mathbb Z})^{\unsure{\prime|*}}$, $\lambda_0\equiv1\ (2)$, tel que $\xi_0=\lambda_0\rho(\lambda_0)^{-1}$ (disons, ou $\xi=\lambda(-l_0^{-1})\rho(\lambda_0)^{-1}$) \deleted{soit tel que} \unsure{satisfasse} dans $\mathrm{Sl}(2,\widehat{\mathbb Z})'$ à une ``équation en $\xi$'' pour $\gamma=\frac{p-1}{2}$ convenable ($p\in\widehat{\mathbb Z}^{*}$) --- i.e. $c^2+d^2+cd\in\widehat{\mathbb Z}^{*}$ (resp. $3(c^2+cd)+d^2\in\widehat{\mathbb Z}^{*}$). Existe-t-il alors un $\unsure{\lambda|\xi}\in\widehat{\Pi}_{0,3}$ au dessus de $\lambda_0$, tel que $\xi=\lambda\rho(\lambda)^{-1}$ (resp.\ldots) \unsure{satisfasse} aussi une équation en $\xi$ (\unsure{nécessairement} avec le même $\gamma$) ?

\scanpage{121}
\mspage{350}

\uline{43} \deleted{Retour sur les notations} -- Digression sur \inserted{les 1-}cocycles des groupes cycliques.

\unsure{Soient|Soit} $\Pi$ un groupe, $\rho$ un \unsure{autom} de $\Pi$ tel que
\[
\rho^n=\mathrm{id} \tag{1}
\]
où $n\in\mathbb N^*$. On a donc une opération du groupe cyclique $\mathbb Z/n\mathbb Z$ sur $\Pi$, d'où un produit semi direct
\[
\widetilde\Pi=\Pi.\,\mathbb Z/n\mathbb Z, \tag{\unsure{1|2}}
\]
et une structure d'extension

\begin{samepage}
(3) \diagram{\begin{tikzcd}[ampersand replacement=\&]
1 \arrow[r] \& \Pi \arrow[r] \& \widetilde\Pi \arrow[r] \& \mathbb Z/n\mathbb Z \arrow[r] \arrow[l,bend left=30,dotted] \& 1.
\end{tikzcd}}
\end{samepage}

On \deleted{identifie}\inserted{note} $\dot\rho$ \deleted{\unsure{avec}}\inserted{le} générateur de $\mathbb Z/n\mathbb Z$, et \unsure{on|à} \inserted{\unsure{on identifie $\rho$ à}}\deleted{\illegible{1}} l'image de celui-ci dans $\widetilde\Pi$, par le scindage canonique \deleted{\illegible{1}}
\[
\mathbb Z/n\mathbb Z\longrightarrow\widetilde\Pi\qquad 1\bmod n\longmapsto\rho. \tag{4}
\]
On a donc
\[
\rho(g)=\rho g\rho^{-1}\qquad(g\in\Pi). \tag{5}
\]
Les scindages de l'extension (3) correspondent biunivoquement aux éléments $\rho'$ de $\widetilde\Pi$ sur $\dot\rho$, tels que
\[
\rho'{}^n=1. \tag{6}
\]
Dire que $\rho'$ est sur $\unsure{\rho|\dot\rho}$ signifie qu'on a
\[
\rho'=\rho'_g\overset{\unsure{\mathrm{dfn}}}{=}g\rho\qquad\text{avec }g\in\Pi, \tag{7}
\]
et en termes d'un tel $g\in\Pi$ l'équation (6) \unsure{Se|se} récrit ainsi
\[
\boxed{g\rho(g)\rho^2(g)\cdots\rho^{n-1}(g)=1} \tag{8}
\]

\scanpage{122}

\deleted{Posons, pour $g\in\Pi$}

\deleted{On aura, pour $\alpha\in\Pi$,}
\deleted{$\unsure{\alpha\rho\alpha^{-1}=}\rho'_g=g\rho$}

\deleted{de sorte que \illegible{2}}

\deleted{On aura, pour \illegible{1} $g$ comme dessus.}

Le groupe $\Pi$ opère par conjugaison sur l'\unsure{ens.} des scindages de l'extension (3), ou ce qui revient au $\overline{\mathrm m}$, sur l'\unsure{ens.} des $\rho'$ au-dessus de $\unsure{\rho|\dot\rho}$, satisfaisant (6), \deleted{On} On aura d'ailleurs, pour $\alpha\in\Pi$
\[
\begin{aligned}
\alpha\rho'_g\alpha^{-1}
&=\deleted{\unsure{(\alpha g\alpha^{-1})(\alpha\rho\alpha^{-1})=\alpha}}\inserted{\alpha g\rho\alpha^{-1}=\alpha g\rho\alpha^{-1}\rho^{-1}\rho}\\
&=(\alpha g\rho(\alpha)^{-1})\rho
=\unsure{\rho'_{\alpha g\rho(\alpha)^{-1}}}
\end{aligned}\tag{9}
\]
i.e. l'opération par conjugaison de $\Pi$ sur l'\unsure{ens.} des scindages de (3), s'interprète comme l'opération $\Theta=\Theta_\rho$ de $\Pi$ sur l'\unsure{ens.} des solutions de (8) \inserted{donnée, pour $\alpha\in\Pi$, par}
\[
\unsure{\Theta_\rho(\alpha)}:g\longmapsto\alpha g\rho(\alpha)^{-1} \tag{10}
\]
(On dira que $\alpha g\rho(\alpha)^{-1}$ est le $\rho$-conjugué de $g$ par l'élément $\alpha\in\Pi$). Les classes de $\Pi$-conjugaison \unsure{de|des} scindages de (3) correspondent donc aux classes de $\Pi$-$\rho$-conjugaison des solutions de (8).

Parmi ces classes, la classe \inserted{``triviale''} \deleted{\unsure{de l'extension}}\inserted{\unsure{du scindage}} canonique de (3), i.e. la classe de $g=1$, est donnée paramétriquement \deleted{par} comme l'\unsure{ens.} des
\[
g=\alpha\rho(\alpha)^{-1}\qquad(\alpha\in\Pi). \tag{11}
\]
Soient $g,g'$ deux solutions \unsure{congruentes|conjuguées} de

\scanpage{123}
\mspage{351}

l'équation (8), de sorte qu'on a
\[
g'=\alpha g\rho(\alpha)^{-1}\qquad\text{i.e. }g'=\Theta(\alpha).g \tag{12}
\]
pour un $\alpha\in\Pi$ convenable. Quand a-t-on unicité \unsure{des|de} $\alpha$ -- plus généralement, que peut-on dire, pour $g,g'$ fixés, de l'\unsure{ens.} des $\alpha\in\Pi$ \deleted{qui} tels qu'on ait (12) (qu'on peut appeler le $\rho$-\uline{transporteur} de $g$ dans $g'$, dans le groupe $\Pi$). \deleted{Comme}\inserted{C'est} évidemment une classe à droite \unsure{modulo|module} le sous-groupe $\Pi(\rho,g)$ de $\Pi$ qui \deleted{stabilise} stabilise $g$
\[
\begin{aligned}
\Pi(\rho,g)&=\{\beta\in\Pi\mid\Theta_\rho(\beta).g=g\}\\
&=\{\beta\in\Pi\mid\beta g\rho(\beta)^{-1}=g\}\\
&=\{\beta\in\Pi\mid\unsure{\{\mathrm{int}(\beta)\}}(g\rho)=g\rho\}\\
&=\Pi^{\deleted{g}\rho'_g}=\{\beta\in\Pi\mid\mathrm{int}(\underset{\inserted{\rho'_g}}{\uline{g\rho}})\beta=\beta\}\\
&=\text{Centralisateur de }\deleted{g}\rho'_g\ \uline{\text{dans }\Pi}.
\end{aligned}\tag{13}
\]

En d'autres termes, si $\alpha$ satisfait (12), l'\unsure{ens.} des $\alpha'$ qui satisfont la $\overline{\mathrm m}$ relation n'est autre que l'\unsure{ens.} $\alpha\Pi(\rho,g)$ :
\[
\begin{aligned}
\alpha'g\rho(\alpha')^{-1}=\alpha g\rho(\alpha)^{-1}
&\Longleftrightarrow\alpha'\in\alpha\Pi(\rho,g)=\alpha\Pi^{\rho'_g}\\
&\text{i.e. }\alpha^{-1}\alpha'\in\Pi^{\rho'_g}
\end{aligned}\tag{14}
\]
Donc on a unicité des $\alpha$ satisfaisant (12) \deleted{\unsure{ssi}} \unsure{ssi}
\[
\Pi^{\rho'_g}=1, \tag{15}
\]
condition qui ne dépend que de la classe de conjugaison du scindage \deleted{\unsure{envisagé}} de (3) défini par $g$ i.e. par $\rho'_g$.

\scanpage{124}

En fait, non seulement le groupe $\Pi$, mais le groupe $\widetilde\Pi$ lui-$\overline{\mathrm m}$ opère par conjugaison sur l'\unsure{ens.} des scindages de l'extension (3), Il en est en particulier \inserted{ainsi} de $\rho$, et on a
\[
\rho(\rho'_g)\overset{\unsure{\mathrm{dfn}}}{=}
\rho\rho'_g\rho^{-1}=\rho(g\rho)\rho^{-1}
=(\rho g\rho^{-1}).\rho=\unsure{g'_{\rho(g)}|\rho'_{\rho(g)}}, \tag{16}
\]
i.e. l'opération de $\rho$ sur l'\unsure{ens.} des scindages en question, s'identifie à l'opération \deleted{de} induite \deleted{par}\unsure{par|$\rho$} la permutation $\rho$ de $\Pi$ sur l'\unsure{ens.} des solutions de (8) $(g\in\Pi)$. Notons d'ailleurs que $g$ et $\rho(g)$ \deleted{\illegible{1}} sont $\Pi$-conjugués, car
\[
\rho(g)=\deleted{g}\inserted{\alpha}g\rho(\alpha)^{-1}\qquad\text{avec }\alpha=g^{-1}. \tag{17}
\]

Notons que \deleted{\illegible{4}}\inserted{l'opération de $\widetilde\Pi$}\deleted{\unsure{$\rho$}}, sur l'\unsure{ens.} des scindages de (3), ou encore sur l'\unsure{ens.} des solutions de (8), est \unsure{tjrs|donc} \unsure{donnée} par
\[
\Theta_\alpha:g\longmapsto\alpha g\rho(\alpha)^{-1} \tag{10}
\]
où maintenant on suppose $\alpha\in\widetilde\Pi$ (et pas nécessairement $\alpha\in\Pi$). Les classes de $\widetilde\Pi$-conjugaison sont \unsure{réduites|identiques} aux classes de $\Pi$-conjugaison, \inserted{(comme il résulte} \unsure{résulte} de (17)) \deleted{\unsure{vient de le voir.}} En particulier, si $\alpha\in\widetilde\Pi$, l'élément
\[
\alpha\rho(\alpha)^{-1}\qquad(\alpha\in\widetilde\Pi) \tag{18}
\]
de $\widetilde\Pi$ est dans $\Pi$, et il \deleted{\illegible{2}} satisfait (8), et

\scanpage{125}
\mspage{352}

se trouve dans la classe principale, i.e. on aura
\deleted{\unsure{(19)}}
\[
\alpha\rho(\alpha)^{-1}=\alpha_0\rho(\alpha_0)^{-1}\qquad\text{pour }\alpha_0\in\Pi\text{ conv.}
\]
En fait, on peut écrire $\alpha=\alpha_0\rho^i$\inserted{, avec $\alpha_0\in\Pi$, $i\in\mathbb Z$}, d'où $\rho(\alpha)^{-1}=\rho^{-i}\rho(\alpha_0)^{-1}$ d'où (19). Plus généralement, on aura, \deleted{\unsure{si $\alpha,\alpha_0$}}
\[
\alpha=\alpha_0\rho^i\qquad(\deleted{\unsure{\alpha\in\widetilde\Pi}}\alpha_0\in\Pi,\ i\in\mathbb Z)
\]
la relation
\deleted{\unsure{(19)}}
\[
\begin{aligned}
\alpha g\rho(\alpha^{-1})
&=\alpha_0\underset{\deleted{\unsure{\rho^i(g)}}}{\underbrace{\rho^i g\rho^{-i}}}\unsure{\rho(\alpha_0^{-1})|\rho(\alpha_0)^{-1}}\\
&=\deleted{\unsure{\alpha_0\rho^i g\rho(\alpha_0\rho^i)^{-1}}}\inserted{\alpha_0\rho^i(g)\rho(\alpha_0)^{-1}}
\end{aligned}
\]
\unsure{avec|car}
\[
\rho^i(g)=\beta g\rho(\beta)^{-1}\qquad
\beta=(g\rho(g)\cdots\rho^{i-1}(g))^{-1}
\]
donc
\[
\left\{\begin{aligned}
\alpha g\rho(\alpha^{-1})&=(\alpha_0\beta)g\rho(\alpha_0\beta)^{-1}\\
\beta&=(g\rho(g)\cdots\rho^{i-1}(g))^{-1}
\end{aligned}\right. \tag{19}
\]
(pour $\alpha=\alpha_0\rho^i$, $i\in\unsure{\mathbb N|\mathbb N^*}$).

Bien sûr, la représentation (18) d'une solution de l'équation (8), \deleted{\unsure{n'est pas}} quand on ne suppose plus $\alpha\in\Pi$ mais seulement $\alpha\in\widetilde\Pi$, n'est plus biunivoque, puisque pour $\alpha=\rho$ p.ex. on trouve \deleted{\unsure{l'élément 1}} $g=\rho\rho^{-1}=1$.

\deleted{Supposons \unsure{de} plus}\inserted{Supposons} plus généralement qu'on dispose d'un groupe
\deleted{\unsure{(20)}}
\[
\unsure{\mathcal E|\varepsilon}\supset\widetilde\Pi
\]
(comme pour le \deleted{\unsure{gr.}} cas où $\Pi=\Pi_{0,3}$, $\widetilde\Pi$ image inverse de $\{1,\dot\rho,\dot\rho^2\}\subset\mathfrak S_3$ dans $\mathfrak S_{0,3}$ extension de $\mathcal T_{0,3}$ par $\Pi_{0,3}$), dans lequel $\Pi$ soit invariant, de sorte que \deleted{$\unsure{\mathcal E|\varepsilon}$ opère par conjugaison}\inserted{l'on a une extension}

\scanpage{126}

\deleted{sur l'\unsure{ens} des $\rho'$ dans $\Pi$ \unsure{tel|tels} que}

\[
(20)\qquad 1\longrightarrow\Pi\longrightarrow\unsure{\mathcal E|\varepsilon}\longrightarrow G\longrightarrow1
\]

\margin{NB l'extension (3) est l'image inv. de (20) par l'inclusion $\mathbb Z/n\mathbb Z\hookrightarrow G$}

\unsure{avec}
\[
(21)\qquad \mathbb Z/n\mathbb Z\hookrightarrow G,\qquad
1\bmod n\longmapsto\dot\rho\in G
\]
\inserted{(élément d'ordre $n$)}
$\widetilde\Pi$ s'identifiant à l'image inv. de $\mathbb Z/n\mathbb Z$ dans $\unsure{\mathcal E|\varepsilon}$.

On ne peut dire en général, si $\rho'$ est un \unsure{élément} de $\unsure{\mathcal E|\varepsilon}$ en dessus de $\dot\rho$ (donc $\rho'\in\widetilde{\unsure{\mathcal E|\varepsilon}}$).
($\overline{\mathrm m}$ s'il satisfait $\rho'{}^n=1$), que pour $\forall$ $\alpha\in\unsure{\mathcal E|\varepsilon}$, $\operatorname{int}(\alpha)(\rho')$ soit encore en dessus de $\dot\rho$ :
pour qu'il en soit ainsi, il \unsure{f. et s.} que $\dot\alpha\in G$ centralise $\dot\rho$. Soit donc $\unsure{\mathcal E|\varepsilon}_\rho$
\[
(22)\qquad
\unsure{\mathcal E|\varepsilon}_\rho=\{\alpha\in\unsure{\mathcal E|\varepsilon}\mid[\dot\alpha,\dot\rho]=1\}
\]
c'est un sous-groupe de $\unsure{\mathcal E|\varepsilon}$ contenant $\widetilde\Pi$, qui
\deleted{\unsure{pour $\forall$}}\inserted{\unsure{lui}} opère sur l'ens des \deleted{\unsure{$\rho'$}} $\rho'$ en dessus de $\dot\rho$,
\inserted{notamment sur \deleted{\illegible{1}} l'ens de \illegible{1} \unsure{qui}}
\unsure{satisfaisant} (8), \deleted{et il} par la $\overline{\mathrm m}$ opération
\[
(23)\qquad
\alpha\rho'_g\alpha^{-1}=\rho'_{\alpha g\rho(\alpha)^{-1}}
\]
de sorte qu'on peut aussi considérer cette \unsure{op.} comme donnée par
\[
(24)\qquad
\left\{
\begin{array}{l}
\Theta_{\rho,\unsure{\mathcal E|\varepsilon}_\rho}(\alpha)=\Theta(\alpha):
g\longmapsto\alpha g\rho(\alpha)^{-1}\\
(\alpha\in\unsure{\mathcal E|\varepsilon}_\rho).
\end{array}
\right.
\]

Cette fois, une opération $\Theta(\alpha)$ pourra permuter les classes de $\Pi$-conjugaison des \deleted{\unsure{div.}} solutions, de (8) dans $\Pi$.
\unsure{On|Il} \deleted{\illegible{1}} aura sur $\unsure{\mathcal E|\varepsilon}_\rho$

\scanpage{127}
\mspage{353}

une structure d'extension
\[
(25)\qquad
\begin{array}{ccccccccc}
1&\longrightarrow&\Pi&\longrightarrow&\unsure{\mathcal E|\varepsilon}_\rho&\longrightarrow&G^\rho&\longrightarrow&1\\
&&&&&&\Vert\\
&&&&&&\multicolumn{3}{l}{\operatorname{Centr}_G(\rho)=\{g\in G\mid\rho g\unsure{g|\rho}^{-1}=g\}}\\
&&&&&&\multicolumn{3}{l}{\phantom{\operatorname{Centr}_G(\rho)}=\{g\in G\mid[\rho,g]=1\}}
\end{array}
\]
et $\unsure{\mathcal E|\varepsilon}_\rho$ opère sur l'ens des classes de $\Pi$-conjugaison de scindages de (3), (ou les classes de $\Pi$-$\rho$-conj. de solutions de (8)) via une opération de $G^\rho$ dessus.

Soit $\deleted{\unsure{g}}\alpha\in\unsure{\mathcal E|\varepsilon}$, considérons l'élément
\[
\deleted{\unsure{g}}\alpha\rho(\deleted{\unsure{g}}\alpha)^{-1}
=\alpha(\rho\alpha\rho^{-1})^{-1}
=\alpha\rho\alpha^{-1}\rho^{-1}=[\alpha,\rho]
\]
Il est dans $\Pi$ ssi on a \deleted{\unsure{que}}
\[
[\dot\alpha,\dot\rho]=1\qquad\text{i.e.}\quad
\dot\alpha\in G^\rho\qquad\text{i.e.}\quad\alpha\in\unsure{\mathcal E|\varepsilon}_\rho.
\]

Pour une solution $g$ de l'équation (8), il pourra arriver qu'on ne puisse \unsure{la|le} mettre sous la forme \inserted{standard} $\alpha\rho(\alpha)^{-1}$ avec $\alpha\in\Pi$, mais qu'on y arrive avec $\alpha\in\unsure{\mathcal E|\varepsilon}_\rho$.
Pour qu'il existe un tel $\alpha$, il f. et s. que la classe de $\Pi\rho$-conj. \inserted{$[g]_\rho$} de $g$ (parmi les solutions de (8)) soit dans l'orbite de la classe de conjugaison triviale \inserted{$[1]_\rho$}, pour l'action considérée de $G^\rho$ sur cet ens de classes de conjugaison ; \deleted{et}\inserted{alors} les $\alpha\in\unsure{\mathcal E|\varepsilon}_\rho$ qui \unsure{satisfont}
\[
(26)\qquad g=\alpha\rho(\alpha)^{-1}
\]

\scanpage{128}

\deleted{sont tels que}
\deleted{satisfont à}
\deleted{\unsure{sont des} \illegible{1}}
\inserted{\unsure{satisfont} à la relation}
\[
(27)\qquad\dot\alpha[1]_\rho=[g]_\rho,
\]
et pour $\forall\ \dot\alpha\in G^\rho$ satisfaisant à cette relation (27), il existe au moins un représentant $\alpha\in\unsure{\mathcal E|\varepsilon}_\rho$ satisfaisant (6) (ce qui est \unsure{pure} tautologie).

A quelle condition a-t-on \uline{unicité} \unsure{du|des} $\alpha\in\unsure{\mathcal E|\varepsilon}_\rho$ satisfaisant à (26) ?
L'ens de ces $\alpha$ est une classe à droite sous le \unsure{ss-groupe} de $\unsure{\mathcal E|\varepsilon}_\rho$ \unsure{formé} des $\alpha$ tels que \deleted{\illegible{1}}
\[
\alpha\rho(\alpha)^{-1}=1\qquad\text{i.e.}\quad
\alpha=\rho(\alpha)\qquad\text{i.e.}\quad[\alpha,\rho]=1,
\]
ce groupe n'est autre que \deleted{$\unsure{\mathcal E|\varepsilon}$}
\[
(28)\qquad\unsure{\mathcal E|\varepsilon}^\rho=\operatorname{Centr}_{\unsure{\mathcal E|\varepsilon}}(\rho)\subset\unsure{\mathcal E|\varepsilon}_\rho.
\]
\unsure{lequel} groupe contient $\rho$ \unsure{donc} le \inserted{\unsure{ss-}}groupe $\simeq\mathbb Z/n\mathbb Z$ \inserted{de $\unsure{\mathcal E|\varepsilon}$} engendré par $\rho$ -- il n'est donc jamais trivial -- i.e. on n'a jamais unicité (sauf dans le cas trivial où \deleted{\unsure{$\rho$}} $n=1$ i.e. $\rho=1$, \deleted{\illegible{1}} mais alors $\unsure{\mathcal E|\varepsilon}^\rho=\unsure{\mathcal E|\varepsilon}$ et il faudrait $\unsure{\mathcal E|\varepsilon}=1$ \dots !)

Mais si on suppose $\Pi^\rho=\{1\}$ i.e.
\[
(29)\qquad\unsure{\mathcal E|\varepsilon}^\rho\cap\Pi=\{1\}
\]
\deleted{\unsure{alors}}
\[
\unsure{\mathcal E|\varepsilon}^\rho_{\deleted{\unsure{\rho}}}
\underset{\uline{\text{injectif}}}{\overset{\mathrm{can}}{\hookrightarrow}}
G^\rho\subset G,
\]
alors les $\alpha\in\unsure{\mathcal E|\varepsilon}^\rho$ solutions de (26) sont en correspon-

\scanpage{129}
\mspage{354}

dance biunivoque avec leurs images \inserted{$\dot\alpha$} dans $G$, i.e. avec le transporteur \inserted{dans $G$} de $[1]_\rho$ \unsure{dans|vers} $[g]_\rho$.
Si on choisit un élément de ce transporteur, on a choisi par \unsure{là-même} un $\alpha$, solution de (26).

\[
\boxed{\begin{gathered}
\text{Dans le cas de }\quad
\Pi=\Pi_{\unsure{0,3}}\inserted{=\Pi_1(U_{0,3})},\qquad
\unsure{\mathcal E|\varepsilon}=\mathfrak S_{0,3}
\inserted{=\Pi_1(U_{0,3},\mathcal T_{0,3})},
\quad\deleted{\unsure{G=\mathcal T_0}}\\
G=\mathcal T_{0,3}\simeq\mathfrak S_3\times\mathbb Z/2\mathbb Z
\end{gathered}}
\]
avec $\rho,\dot\rho$ comme d'accoutumée, l'ens des classes de conjugaison des \deleted{\unsure{solutions}}\inserted{scindages} de (3) i.e. des solutions de (8), \unsure{s'identifie} à $U_{0,3}^{\dot\rho}$, qui a deux éléments, \unsure{les} deux "pôles" $-\overline j,-j$ -- dont l'un, choisi comme point base, correspond à la classe triviale.

\deleted{\unsure{Ici le} \illegible{1} $G^\rho\simeq\unsure{\mathbb Z/2\mathbb Z}\{1,\tau\}\simeq\mathbb Z$ est le groupe}
\inserted{Ici}
\[
G^\rho=\unsure{\mathcal T|\mathcal G}_{0,3}^{\rho}
\simeq\mathbb Z/3\mathbb Z\times\mathbb Z/2\mathbb Z
\]
est le groupe engendré par $\rho$ et $\tau$, et $\rho$ opère \unsure{trivialmt|trivialement} sur l'ens des deux classes, tandis que $\tau$ les permute (i.e. $\mathbb Z/\unsure{2|n}\mathbb Z$ \inserted{engendré par $\tau$} y opère de façon simplement transitive).
Donc si on \unsure{pose}
\[
\begin{aligned}
(30)\qquad\unsure{\mathcal E|\varepsilon}'
&=\inserted{\unsure{\widetilde\Pi|\Pi}_{0,3}=}\Pi\cdot\{1,\tau_\infty\}\\
&=\text{image inverse de }\{1,\tau_\infty\}\subset G\\
&\phantom{={}}\text{dans }\unsure{\mathcal E|\varepsilon}=\mathfrak S_{0,3}
\end{aligned}
\]
\margin{(\deleted{\illegible{1}} \unsure{gr.} cartographique triangulé \unsure{pondéré} non orienté \dots}
on voit que
\[
(31)\qquad\unsure{\mathcal E|\varepsilon}'{}^\rho=\{1\},
\]
\deleted{(conséquence immédiate de la relation $\Pi^\rho=(1)$)}
et l'ens des solutions dans \uline{$\Pi_{0,3}$} de l'équation
\[
(32)\qquad g\rho(g)\rho^2(g)=1
\]

\scanpage{130}

est en corr. 1-1 avec l'ens des éléments de $\widetilde\Pi_{0,3}=\unsure{\mathcal E|\varepsilon}'$ (le groupe cartographique \inserted{triangulé} pondéré non orienté, ayant comme générateurs les éléments $\tau_0,\tau_1,\tau_\infty$ reliés par les relations $\tau_0^2=\tau_1^2=\tau_\infty^2=1$ sans plus, avec $l_0=\tau_\infty\tau_1$, $l_1=\tau_0\tau_\infty$, $l_\infty=\tau_1\tau_0$ \dots)
Il doit encore en être de $\overline{\mathrm m}$ quand on se place dans les groupes $\Pi_{0,3}^{\wedge}$ et $\widetilde\Pi_{0,3}^{\wedge}=\widetilde\Pi_{0,3}\cdot\{1,\tau_\infty\}$, si on admet (chose à vérifier ultérieurement)
\[
(33)\qquad(\widetilde\Pi_{0,3}^{\wedge})^\rho=\{1\}
\]
\inserted{et que les éléments d'ordre fini de $\widehat{\mathfrak S}_{0,3}$ sont conjugués des éléments d'ordre fini de $\mathfrak S_{0,3}$ \dots}

Notons que l'on a ici
\[
(34)\qquad\tau_\infty\rho(\tau_\infty)^{-1}
=\tau_\infty\rho(\tau_\infty)=l_1^{-1},
\qquad\deleted{\unsure{\text{en}}}
\]
en vertu des formules (23) du \S\ 42
\[
\rho_\infty\overset{\unsure{\mathrm{dfn}|\text{déf}}}{=}
\tau_\infty\rho\tau_\infty^{-1}
=l_1^{-1}\rho=\rho l_0^{-1}.
\]

La situation est encore plus simple pour l'\unsure{élémt|élément} $\sigma_\infty$ de $\mathfrak S_{0,3}$, opérant sur $\Pi$ par automorphisme d'ordre $n=2$, car ici toute solution de l'équation \unsure{correspondante}
\[
(35)\qquad\deleted{\unsure{(36)}}\quad
\deleted{\unsure{g}}h\,\sigma_\infty(\deleted{\unsure{g}}h)=1
\qquad(h\in\Pi)
\]
se met sous la forme standard
\[
(36)\qquad h=\beta\sigma_\infty(\beta)^{-1}
\qquad(\beta\in\Pi)
\]
et ceci de \uline{façon unique}. Cela \unsure{résulte} en effet de la théorie générale \dots, compte tenu que
\[
(37)\qquad U_{0,3}^{\dot\sigma_\infty}=\{\unsure{\tfrac12}\}
\]
